\section{Introduction}\label{sec:intro}

An employee at a large company can appear with different colleagues in two register-derived networks even when the underlying personnel records are identical. The original Dutch rule favours colleagues who live nearby. The Danish description uses a random sample and makes sampled ties two-way. A researcher comparing work-layer degree must know whether a limit applies before or after this step. An employee can receive more nominations than they make. This detail changes both degree and the composition of the group whose clustering is measured \citep{vanderlaan2022,cremers2025}.

The mechanism begins with a membership, such as a workplace identifier, household address or school cohort. Each identifies a set of candidate alters. A builder can link the whole group, select some alters or first divide the group by geography or time. If nominations become undirected edges, an incoming nomination can add a tie the recipient did not choose. These operations change edge count, topology and layer overlap independently of any change in the records. Consequently a valid comparison needs both a population input and an explicit rule for deriving ties from it.

Published releases describe their own choices. Several groups use the resulting networks to analyse inequality, cohesion and labour mobility \citep{kazmina2024,bokanyi2026,lengyel2017}. Earlier work also varies specific choices within a release or models the affiliations that give rise to a network \citep{lengyel2017,hamilton2026}. A practical gap remains for an analyst who wants to hold people and memberships fixed and rerun several published constructions, including disputed readings. The implementation must preserve identifiers across outputs and show where a published rule has been inferred or simplified.

builtnet makes construction choices in register-based population networks inspectable and testable by encoding source-attributed alternatives as executable cards that can be applied to the same membership records. It contributes (1) a catalogue covering the original and yearly Dutch releases and the Danish and Swedish constructions, with runnable alternatives, (2) a common table interface and sparse graph implementation with explicit output semantics, and (3) a combination of regression checks, worked examples and edge-free summaries that reveals where competing rules have different implications.

\textbf{Claim boundary.} We claim repeatable comparisons of stated rules on supplied memberships within the supported interface. Reconstructing the national releases, identifying a tie as an observed contact, estimating a causal effect, and guaranteeing nationwide computational performance are outside that claim. The article situates the cards against related systems, describes their operators, examines the existing validation and worked examples, and ends with the conditions under which the comparisons can be reused.
